\section{Smart Contract Design}
\label{sec:contracts}

\subsection{Architecture}

PriceOracle-AVS reuses the two-contract pattern of the Incredible
Squaring template~\cite{layrlabs_isavs}: a
\texttt{PriceOracleServiceManager}, which registers the AVS with the
EigenLayer core contracts and gates operator opt-in, and a
\texttt{PriceOracleTaskManager}, which schedules price-quote tasks
and verifies the BLS-aggregated responses returned by the
Aggregator.  The split keeps the EigenLayer-facing administrative
surface in a small, audit-friendly contract while the
application-specific task logic lives in a separate, easily upgradeable
contract.  Operators register once with the ServiceManager, after
which the TaskManager drives the per-task lifecycle:
\texttt{NewTaskCreated} $\rightarrow$ off-chain quote retrieval and
BLS signing $\rightarrow$ \texttt{respondToTask} $\rightarrow$
optional \texttt{raiseAndResolveDeviationChallenge}.

\subsection{Modified Task and TaskResponse Structs}

The principal change relative to the template is the redefinition of
the \texttt{Task} and \texttt{TaskResponse} structs (Listing~\ref{lst:taskstruct}).
The template's \texttt{numberToBeSquared} field becomes an
\texttt{assetPair} string identifying the feed (\texttt{"ETH/USD"} or
\texttt{"BTC/USD"} in our deployment).  The response now carries the
median price and its observed dispersion rather than a single
squared integer.  Both fields use the 6-decimal fixed-point
convention common to USD pairs in the Chainlink and Pyth
ecosystems~\cite{breidenbach2021chainlink, pyth2023}.

\begin{lstlisting}[language=Solidity, caption={Task and TaskResponse structs.}, label={lst:taskstruct}]
struct Task {
    string   assetPair;       // e.g. "ETH/USD"
    uint32   taskCreatedBlock;
    bytes    quorumNumbers;
    uint32   quorumThresholdPercentage;
}

struct TaskResponse {
    uint32   referenceTaskIndex;
    uint256  priceMedian;     // 6-decimal fixed
    uint256  priceVariance;
}
\end{lstlisting}

\subsection{Deviation Slashing Rule}

The most consequential addition is the deviation slashing function.
For each task index $T$, after the Aggregator submits the consensus
\texttt{priceMedian} $m$, an external Challenger may invoke
\texttt{raiseAndResolveDeviationChallenge}$(T, o_i, p_i)$ to assert
that operator $o_i$ reported a price $p_i$ outside the tolerance band:
\begin{equation}
\text{slash}(o_i) \;\iff\; \frac{|p_i - m|}{m} \;>\; \tau,
\label{eq:slash}
\end{equation}
where $\tau$ is a deployment-time constant fixed at $5\%$
(see Section~\ref{sec:agg-tolerance} for the empirical motivation).
If the inequality holds, the contract calls into the EigenLayer
\texttt{Slasher} to forfeit a configurable fraction of $o_i$'s
restaked balance.  If it does not, the challenger pays a
predetermined bond to disincentivise nuisance challenges, mirroring
the optimistic-oracle pattern of UMA~\cite{uma2020}.  The full
contract source is available in our open-source
fork.\footnote{\url{https://github.com/<TODO M1: paste fork URL>}}

\subsection{Foundry Test Coverage}

We exercise the contracts with a Foundry test
suite~\cite{foundry} of $\geq 10$ unit tests, covering at minimum:
(i) \texttt{NewTaskCreated} event emission for valid asset pairs,
(ii) rejection of below-quorum responses,
(iii) successful BLS aggregate verification,
(iv) slashing triggered when $|p_i - m| / m > \tau$,
(v) slashing rejected (and challenger bond forfeited) when
$|p_i - m| / m \leq \tau$,
and (vi) replay protection against duplicate \texttt{respondToTask}
calls.  Total branch coverage on the modified contracts is
\textit{TBD\% [M1: fill via \texttt{forge coverage}]}; gas
consumption per task lifecycle is \textit{TBD gas [M1: fill via
\texttt{forge test --gas-report}]}.

\begin{figure}[t]
  \centering
  % \includegraphics[width=0.95\columnwidth]{figures/fig_contracts.pdf}
  \framebox[0.95\columnwidth][c]{\raisebox{1.2cm}{\textit{[placeholder: contract architecture --- M1 to draw]}}}
  \caption{PriceOracle-AVS smart-contract architecture.  Operators
           opt in via the ServiceManager; the TaskManager emits price
           tasks, verifies the aggregator's BLS-signed median
           response, and routes outlier challenges to the EigenLayer
           Slasher.}
  \label{fig:contracts}
\end{figure}
